\documentclass{article}

\begin{document}

This short document describes the mathematics that the source code currently
implements. This code effectively implements a gaussian transform from some
gaussian coefficients to real space grid and vis versa. The two transforms are
described by
\begin{equation}
\label{eq:integrate}
V^{c}_{\alpha\beta\gamma} = \int V({\bf r}) (x - x_{c})^\alpha(y - y_{c})^\beta(z - z_{c})^\gamma e^{-\mu_{c} ({\bf r} - {\bf r}_{c})^2},
\end{equation}
and
\begin{equation}
\label{eq:collocate}
V({\bf r}) = \sum_{ij} \sum_{\alpha\beta\gamma} (x - x_{c})^\alpha(y - y_{c})^\beta(z - z_{c})^\gamma e^{-\mu_{c} ({\bf r} - {\bf r}_{c})^2},
\end{equation}
where ${\bf r}_{c}$ is the center of the gaussian, $\alpha,\beta,\gamma$ are
positive integers with $\alpha+\beta+\gamma \leq l$ and $l$ is an integer
describing the angular momentum.

Eq.\ref{eq:collocate} is implemented in the
$\mathtt{collocate\_original\_standalone.cu}$ and Eq.\ref{eq:integrate} in the
$\mathtt{integrate\_original\_standalone.cu}$. The original code implements the
formulas directly and distinguish between rectangular and generic lattices
(orthogonal lattice vectors vs generic lattice vectors) but we know it is
possible to implement these formulas using some tensor contraction or
matrix-matrix multiplications.

A second step is required to transform the coefficients
$V^{ij}_{\alpha\beta\gamma}$ back to the form used by cp2k for the computations.

The integral Eq.\ref{eq:integrate} converges exponentially and is effectively
restricted to an integral over a sphere of given radius to reduce computation
time. It is also possible to exploit the lattice properties such as
orthorhombicity. When the lattice vectors are orthogonal, the integral can be
expressed as a tensor reduction of the form
\begin{equation}
  \label{decomposition::tensor}
  V_{\alpha\beta\gamma} = \sum_{ijk} P^{i}_{\alpha} P^{j}_{\beta} P^{k}_{\gamma} V_{ijk} E_{ijk}
\end{equation}
where $P^{i}_{\alpha} = (x_i - x_c)^\alpha$, $P^{j}_{\beta} = (y_j -
y_c)^\beta$, and $P^{k}_{\gamma} = (z_k - z_c)^\gamma$ are matrices with
$x_i,y_j,z_k$ discrete lattice coordinates and $x_c,y_c, z_c$ the
coordinates of the gaussian center. The tensor $E_{ijk}$ is given by
$$E_{ijk} = \exp\left[ - ({\bf r}_{ijk} - {\bf r}_c)^2\right].$$

We can use the properties of the exponential $$\exp(-\left( a-b \right)^2) =
\exp\left( -a^2\right) \exp\left( -b^2\right) \exp \left( 2 ab \right)$$ to
calculate $E_{ijk}$ efficiently provided that we express all coordinates in the
lattice vector basis instead of the cartesian basis. In both cases, the number
of exponentials to be computed is $n + m + o$ where $n,m,o$ are the dimensions
of the discrete cude containing the ball of radius $r_{\textrm{cut}}$ for the
orthorhombic case and $n (m + o) + m o$ for the general cases. It is
particularly important when the implementation runs on CPU less so on GPU.

the goal is to optimize the implementations of Eq.\ref{eq:integrate} and
Eq.\ref{eq:collocate}.



\end{document}
